% !TeX root = ../常微分方程笔记.tex
% 由主文档通过 \input 载入；请编译主文档。
\setcounter{chapter}{2} % 保留原书章号。
\chapter{方法索引与习题提示}
\label{app:ode-methods}

\section{按方程结构查找方法}
\begin{description}
 \item[可分离与一阶线性] 用式\eqref{eq:ode-separable}或\eqref{eq:ode-first-linear}；前者检查零因子对应的平衡解，后者检查系数连续的区间。
 \item[可代换的一阶方程] 比值型取 $y/t$，Bernoulli 型取 $y^{1-\alpha}$，已知特解的 Riccati 型取 $y-y_p=1/u$；记录每个代换的可逆范围。见第\ref{chap:ode-first-order}章。
 \item[隐式与可降阶方程] 先检查隐函数条件，再考虑导数参数化、缺项代换与能量积分。见第\ref{chap:ode-reduction}章。
 \item[存在、唯一与延拓] 以定理\ref{thm:ode-picard}判断局部问题，用定理\ref{thm:ode-continuation}和先验界处理最大区间，不从光滑性直接推断全局存在。
 \item[高阶线性与线性系统] 常系数用特征根或矩阵指数；非齐次用式\eqref{eq:ode-system-variation}；变系数一般没有简单指数公式。见第\ref{chap:ode-linear-scalar}、\ref{chap:ode-linear-system}章。
 \item[非线性长期行为] 先找平衡点、不变区域与守恒量，再用线性化或 Lyapunov 方法。零实部特征值通常需要额外分析。见第\ref{chap:ode-phase-plane}、\ref{chap:ode-stability}章。
 \item[级数、变换与边值] 常点与正则奇点选择不同级数形式；分段输入可用拉普拉斯变换；边值问题先检查相容性。见第\ref{chap:ode-series}至第\ref{chap:ode-boundary}章。
 \item[数值求解] 分别检查精度、稳定区域与误差控制；通过步长变化、已知解或结构量交叉核验。见第\ref{chap:ode-numerical}章。
\end{description}

\section{完成一道题后的自检}
一个完整解答应明确未知函数、自变量、参数与初始或边界数据；说明关键代换及除法是否合法；将结果代回原方程；给出相关定义域或最大区间；说明定理的假设如何满足。稳定性结论要写清局部或全局、稳定或渐近稳定。数值题还应给出算法、步长或容差、误差度量和参考结果的来源。

\section{部分习题答案与提示}
以下只给出检查结果与关键步骤，计算和证明仍应按正文方法完成。
\begin{description}
 \item[第\ref{chap:ode-intro}章] 习题\ref{ex:ode-classify}依次为二阶线性非自治、一阶非线性自治、一阶非线性非自治。习题\ref{ex:ode-square}：$a\ne0$ 时 $y=a/(1-at)$；$a>0$ 的最大区间为 $(-\infty,1/a)$，$a<0$ 时为 $(1/a,\infty)$；$a=0$ 时为全实轴上的零解。习题\ref{ex:ode-boundary-intro}的解依次为 $1+2t$ 与 $1+t$。
 \item[第\ref{chap:ode-first-order}章] 习题\ref{ex:ode-separate}为 $3e^{t^2}$；习题\ref{ex:ode-linear}为 $(e^t-e^{-2t})/3$。习题\ref{ex:ode-exact}可取势函数 $H=t^2y+t+y^2$，当 $t^2+2y\ne0$ 时可局部解出 $y$。
 \item[第\ref{chap:ode-reduction}章] 习题\ref{ex:ode-clairaut}的直线族为 $y=Ct-C^2$，奇解为 $t^2/4$；习题\ref{ex:ode-reduce}为 $(e^{2t}-1)/2$。
 \item[第\ref{chap:ode-existence}章] 习题\ref{ex:ode-picard}对递推积分作归纳即可，每次积分增加下一项 $t^{k+1}/(k+1)!$。
 \item[第\ref{chap:ode-dependence}章] 习题\ref{ex:ode-growth}中 $|f|\leq1$、状态 Lipschitz 常数可取 $1$，故初值差的上界为 $e^{|t-t_0|}$ 倍原差。习题\ref{ex:ode-domain}为 $\sqrt{1+2t}$，最大区间 $(-1/2,\infty)$；左端趋近右端函数未定义的 $y=0$。
 \item[第\ref{chap:ode-linear-scalar}章] 习题\ref{ex:ode-constant}为 $\cos2t+\sin2t$；习题\ref{ex:ode-resonance}可取 $-t\cos t/2$，原试探式属于齐次解空间。
 \item[第\ref{chap:ode-linear-system}章] 习题\ref{ex:ode-system}为 $x_1=(e^{2t}+1)/2$、$x_2=(e^{2t}-1)/2$；习题\ref{ex:ode-nilpotent}为 $I+tA+t^2A^2/2$。
 \item[第\ref{chap:ode-phase-plane}章] 习题\ref{ex:ode-phase}中 $0,2$ 吸引，$1$ 排斥；习题\ref{ex:ode-spiral}为逆时针旋转的吸引焦点，可代入 $(1,0)$ 检查向上运动。
 \item[第\ref{chap:ode-stability}章] 习题\ref{ex:ode-stable}依次为稳定但非渐近稳定、指数稳定、渐近稳定但非指数稳定、不稳定。习题\ref{ex:ode-lyapunov}有 $\dot V=-u^2-v^2$，且 $V$ 径向无界。
 \item[第\ref{chap:ode-series}章] 习题\ref{ex:ode-airy}为 $1+t^3/6+t^6/180+t^9/12960$；习题\ref{ex:ode-frobenius}的指标根为 $\pm2$，通解 $C_1t^2+C_2t^{-2}$。
 \item[第\ref{chap:ode-laplace}章] 习题\ref{ex:ode-laplace}为 $(1-e^{-2t})/2$；习题\ref{ex:ode-delay}为 $2e^{-2s}/s^3$，$s>0$。
 \item[第\ref{chap:ode-boundary}章] 习题\ref{ex:ode-neumann}为 $\lambda_n=n^2$、$y_n=\cos nt$，$n=0,1,\ldots$；习题\ref{ex:ode-mixed}为 $\lambda_n=(n+1/2)^2$，$n=0,1,\ldots$。
 \item[第\ref{chap:ode-numerical}章] 习题\ref{ex:ode-euler-value}为 $1.1^{10}\approx2.593742$，小于 $e$；习题\ref{ex:ode-euler-stability}要求 $0<h<0.1$，两给定步长的放大因子分别为 $0$ 与 $-2$。
 \item[第\ref{chap:ode-modeling}章] 习题\ref{ex:ode-harvest}可用 $u'=1/4-h-(u-1/2)^2$ 分类，再参照本章的相线分析；习题\ref{ex:ode-sir}中 $\mathcal R_0S(0)/N=0.8<1$，故初期感染人数减少。
\end{description}

